\documentclass[11pt]{article}
\usepackage[margin=0.85in]{geometry}
\usepackage{amsmath,amssymb}

\title{Disproof of Conjecture 00000002731\\
\large Equal \(h\)-vectors and different simplicial surjection behavior}
\author{}
\date{}

\begin{document}
\maketitle

Consider the following two one-dimensional simplicial complexes on the
vertices \(0,1,2,3\). Their maximal faces are the listed edges.
\[
\begin{array}{c|l|c}
\text{Complex} & \text{Edges} & \text{Chromatic number}\\ \hline
C & 01,\ 12,\ 23,\ 03 & 2\\
T & 01,\ 02,\ 12,\ 23 & 3
\end{array}
\]
Both are connected and pure because they contain the spanning path
\(0\!-\!1\!-\!2\!-\!3\). The coloring with vertex colors \(0,1,0,1\)
is proper for \(C\), and its edge \(01\) requires at least two colors.
The triangle \(012\) in \(T\) requires three colors, and the assignment
\(0,1,2,0\) is a proper three-coloring.

\section*{Equal \(h\)-vectors}

For either complex \(K\) we have \(f_{-1}=1\), \(f_0=4\), and \(f_1=4\).
The defining identity for its \(h\)-vector gives
\[
  \sum_{i=0}^{2} h_i t^{2-i}
   = \sum_{i=0}^{2} f_{i-1}(t-1)^{2-i}
   = (t-1)^2+4(t-1)+4 = t^2+2t+1 .
\]
Thus both full \(h\)-vectors are \((1,2,1)\).

\section*{Different surjection behavior}

A simplicial map to a simplex is dimension-preserving exactly when
the endpoints of each source edge have distinct images. For a source
with an edge, the target therefore needs at least two vertices.
Surjectivity on faces requires a source face to map to the full target
simplex, which allows at most two target vertices. For surjectivity on
geometric realizations, the same upper bound follows because the target
barycenter lies in the image of a source face only if that face maps to
all target vertices. In either sense the target must be a \(1\)-simplex.

For this target, a dimension-preserving map is a proper two-coloring.
The image of any source edge is the full target edge, so such a map
is surjective. The coloring of \(C\) gives a dimension-preserving
surjection, while the triangle in \(T\) rules one out. Thus the common
\(h\)-vector cannot determine existence of the surjection or the
chromatic number. In particular any ascending-segment statistic of
that vector agrees for the two complexes and cannot give the stated
criterion.

\section*{Lean verification}

The Lean~4.19.0 project lists the faces via their vertices and edges.
It proves the target-size restriction, computes both \(h\)-vectors,
establishes the two chromatic numbers, and rules out any criterion
depending only on the \(h\)-vector. Run \verb|lake build| in
\verb|lean/| to verify the proofs and print their axiom dependencies.

\end{document}
